\section{Local response hierarchy}
\label{sec:localresults}

The first central test compares the EOS response transverse to the stellar sequence in the I--Love and C--Love planes. Both are constructed from the same stellar configurations and EOS metric, and no empirical relation is used to define either normal direction.

For each reference configuration we compute the covariance-weighted transverse response and record its nonzero normal eigenvalue in the two-dimensional plane. The key diagnostic is therefore not the raw scatter of a sampled EOS catalog, but the ratio of local EOS response scales for I--Love and C--Love at matched stellar state.

\todo{Stage 3: report the transverse response as a function of mass or central enthalpy, including variation under EOS metric, observable metric, basis size, and numerical tolerance.}

A decisive positive result would be a robust hierarchy in which the normal EOS response of I--Love is substantially smaller than that of C--Love across the ordinary neutron-star domain. A null result would imply that the observed global universality is not captured by the proposed local metric construction, or that the chosen EOS/observable metrics do not represent the relevant physical perturbations.

\section{Integrability and global reconstruction}
\label{sec:global}

The low-response covector field is transported across the stellar domain and integrated without reference to the standard I--Love fitting formula. In the two-dimensional I--Love plane the local normal field is integrable, but a nontrivial test is obtained by embedding the construction in the full \((\ln C,\ln\Ibar,\ln\Lamtwo)\) space and checking path dependence after quotienting the stellar-sequence direction.

\todo{Stage 4: add integrability diagnostic, path-dependence statistic, and reconstructed I--Love curve. Compare to the empirical relation only after reconstruction.}

\section{Curvature and finite EOS scatter}
\label{sec:curvature}

The second-order tensor entering Eq.~\eqref{eq:second-order} is evaluated using directional Hessian-vector products. Its covariance-weighted contraction gives a local prediction for finite EOS scatter after first-order cancellation.

\todo{Stage 5: compare curvature-predicted scatter with direct nonlinear EOS draws over increasing perturbation amplitude.}

\section{Controlled breakdown}
\label{sec:breakdown}

Flexible EOS models and sharp phase-transition-like structure provide a stress test rather than an extension of the universality claim. We will determine whether loss of I--Love accuracy is accompanied by a growing transverse response, loss of integrability, enhanced curvature, or some combination of these effects. The density localization of the functional kernels will be used to identify which EOS regions produce the breakdown.

\todo{Stage 6: insert broad-EOS and phase-transition results, with the same diagnostics used in the baseline analysis.}
